% Solutions/hints for ch18 exercises (assembled into Appendix C)

\begin{solution}{exr:ch18-hand}
Predict: $\hat{\state}^-_2=\mat{F}\hat{\state}_1=(1.195+0.859,\ 0.276+0.143,\ 0.859,\ 0.143)=(2.054,0.419,0.859,0.143)$; $a^-=0.2226+2\cdot0.1150+0.6172+0.0333=1.1031$, $b^-=0.1150+0.6172+0.05=0.7822$, $c^-=0.6172+0.1=0.7172$. Update: $S=1.1031+0.25=1.3531$, $k_p=1.1031/1.3531=0.8152$, $k_v=0.7822/1.3531=0.5781$; $\vect{y}_2=(1.71-2.054,\ 0.92-0.419)=(-0.344,0.501)$; $\hat{\state}_2=(2.054-0.8152\cdot0.344,\ 0.419+0.8152\cdot0.501,\ 0.859-0.5781\cdot0.344,\ 0.143+0.5781\cdot0.501)=(1.774,0.827,0.660,0.433)$; $a=(1-0.8152)\cdot1.1031=0.2038$, $b=0.1848\cdot0.7822=0.1445$, $c=0.7172-0.5781\cdot0.7822=0.2650$. All agree with \cref{tab:ch18-trace} to the rounding of the hand computation. Step~3 prediction: $\hat{\state}^-_3=(1.774+0.660,\ 0.827+0.433,\ 0.660,\ 0.433)=(2.434,1.260,0.660,0.433)$. Its $x$-component is below the true $3$ because the measurement $\meas_2$ was $0.29$~m short of the true position and the filter, still trusting the sensor heavily ($k_p=0.82$), followed it, and because the velocity estimate $0.66$~m/s is still well below the true $1$~m/s: the velocity is only seen through position differences and needs several steps to converge (it is $1.13$ after step~5).
\end{solution}

\begin{solution}{exr:ch18-gain}
(a) $\mat{P}(\Kgain)=\mat{P}^--\Kgain\mat{H}\mat{P}^--\mat{P}^-\mat{H}\T\Kgain\T+\Kgain\mat{S}\Kgain\T$ with $\mat{S}=\mat{H}\mat{P}^-\mat{H}\T+\mat{R}$. Taking the trace and differentiating with the two rules, $\partial\tr\mat{P}(\Kgain)/\partial\Kgain=-2\mat{P}^-\mat{H}\T+2\Kgain\mat{S}$, which vanishes for $\Kgain^*=\mat{P}^-\mat{H}\T\mat{S}\inv$. (b) By \cref{eq:ch18-best-gain}, $\mat{P}(\Kgain)-\mat{P}(\Kgain^*)=(\Kgain-\Kgain^*)\mat{S}(\Kgain-\Kgain^*)\T$ is positive semidefinite for every $\Kgain$, so every diagonal entry of $\mat{P}(\Kgain)$ is at least the corresponding entry of $\mat{P}(\Kgain^*)$, and the trace is minimised; a stationary point of a function bounded below by its value there is a global minimum. (c) At the optimum $\Kgain^*\mat{S}=\mat{P}^-\mat{H}\T$, so $\Kgain^*\mat{S}\Kgain^{*\mathsf{T}}=\mat{P}^-\mat{H}\T\Kgain^{*\mathsf{T}}=(\Kgain^*\mat{H}\mat{P}^-)\T=\Kgain^*\mat{H}\mat{P}^-$, the last step because $\Kgain^*\mat{H}\mat{P}^-=\mat{P}^-\mat{H}\T\mat{S}\inv\mat{H}\mat{P}^-$ is symmetric. Substituting into the expansion of (a), the two middle terms and the last combine to $-\Kgain^*\mat{H}\mat{P}^-$, giving $(\mat{I}-\Kgain^*\mat{H})\mat{P}^-$, and $\mat{P}^--\Kgain^*\mat{S}\Kgain^{*\mathsf{T}}$ is the same matrix.
\end{solution}

\begin{solution}{exr:ch18-q}
(a) A unit impulse of jerk at time $\tau$ raises the acceleration by $1$ from $\tau$ on, hence the velocity at $\dt$ by $\dt-\tau$ and the position by $(\dt-\tau)^2/2$; superposing gives the stated integral. With $s=\dt-\tau$ and $\E[j(\tau)j(\tau')]=q\,\delta(\tau-\tau')$, $\mat{Q}=q\int_0^{\dt}(s^2/2,\,s,\,1)\T(s^2/2,\,s,\,1)\,ds$, whose entries are $q\int_0^{\dt}s^4/4=q\dt^5/20$, $q\int s^3/2=q\dt^4/8$, $q\int s^2/2=q\dt^3/6$, $q\int s^2=q\dt^3/3$, $q\int s=q\dt^2/2$ and $q\dt$. (b) Direct multiplication: with $\mat{F}=\begin{psmallmatrix}1&\dt\\0&1\end{psmallmatrix}$ and $\mat{Q}(\dt)=q\begin{psmallmatrix}\dt^3/3&\dt^2/2\\ \dt^2/2&\dt\end{psmallmatrix}$, $\mat{F}\mat{Q}\mat{F}\T=q\begin{psmallmatrix}\dt^3/3+\dt^3+\dt^3&\dt^2/2+\dt^2\\ \dt^2/2+\dt^2&\dt\end{psmallmatrix}$, and adding $\mat{Q}(\dt)$ gives $q\begin{psmallmatrix}8\dt^3/3&2\dt^2\\2\dt^2&2\dt\end{psmallmatrix}=\mat{Q}(2\dt)$. For the induction, note that the integral form $\mat{Q}(T)=\int_0^T e^{\mat{A}s}\mat{G}q\mat{G}\T e^{\mat{A}\T s}ds$ splits at $s=T-\dt$ into $\mat{Q}(T-\dt)$ transported by $\mat{F}(\dt)$ plus $\mat{Q}(\dt)$; equivalently, the sum over $i<j$ is a Riemann sum of the same integral. This is why $q$ is the invariant parameter and $\mat{Q}$ is not. (c) \code{process\_noise\_numeric(A, [0,0,1], q, dt)} with the $3\times3$ nilpotent $\mat{A}$ matches \code{ca\_model} to $10^{-6}$; the self-test of \code{ch18\_kalman.py} contains both checks.
\end{solution}

\begin{solution}{exr:ch18-init}
(a) $\hat{\state}_0$ is the linear map $\begin{psmallmatrix}\mat{0}&\mat{I}\\ -\mat{I}/\dt&\mat{I}/\dt\end{psmallmatrix}$ of $(\meas_a,\meas_b)$, whose covariance is $\diag(\mat{R},\mat{R})$; the product $\mat{M}\diag(\mat{R},\mat{R})\mat{M}\T$ has blocks $\mat{R}$, $\mat{R}/\dt$, $\mat{R}/\dt$ and $(\mat{R}+\mat{R})/\dt^2$. (b) Per axis, with $q=0$: predicting from $(a,b,c)=(r,0,V)$ gives $a^-=r+\dt^2V$, $b^-=\dt V$, $c^-=V$, hence $S=2r+\dt^2V$, $k_p=(r+\dt^2V)/S\to1$ and $k_v=\dt V/S\to1/\dt$ as $V\to\infty$. The updated mean is $\hat p=z_a+k_p(z_b-z_a)\to z_b$ and $\hat v=k_v(z_b-z_a)\to(z_b-z_a)/\dt$. The updated covariance is $a=(1-k_p)a^-=r(r+\dt^2V)/(2r+\dt^2V)\to r$, $b=(1-k_p)b^-=r\dt V/(2r+\dt^2V)\to r/\dt$ and $c=c^--k_vb^-=V-\dt^2V^2/(2r+\dt^2V)=2rV/(2r+\dt^2V)\to2r/\dt^2$, which is \cref{eq:ch18-init2}. (c) Starting from $\meas_b$ with a zero velocity of variance $v_{\max}^2$ discards the information in $\meas_a$ and, unless $v_{\max}$ is large compared with the true speed, biases the first velocity estimates towards zero; the two-point start has a much larger initial velocity variance ($2\sigma_p^2/\dt^2$) but is unbiased and correctly correlated with the position, and in experiment A its velocity error falls below $1$~m/s within about two seconds, while the zero-velocity start with $v_{\max}=5$~m/s takes noticeably longer.
\end{solution}

\begin{solution}{exr:ch18-scales}
(a) $\operatorname{Var}(\lambda z_1+(1-\lambda)z_2)=\lambda^2\sigma_1^2+(1-\lambda)^2\sigma_2^2$; setting the derivative $2\lambda\sigma_1^2-2(1-\lambda)\sigma_2^2$ to zero gives $\lambda^*=\sigma_2^2/(\sigma_1^2+\sigma_2^2)$ and the minimum $\sigma^2=\sigma_1^2\sigma_2^2/(\sigma_1^2+\sigma_2^2)$, that is $1/\sigma^2=1/\sigma_1^2+1/\sigma_2^2$: precisions add. (b) Write $K=1-\lambda^*=\sigma_1^2/(\sigma_1^2+\sigma_2^2)$; with $z_1=\hat x^-$, $\sigma_1^2=P^-$, $z_2=z$ and $\sigma_2^2=R$ this is $K=P^-/(P^-+R)$ and the estimate is $(1-K)\hat x^-+Kz=\hat x^-+K(z-\hat x^-)$, exactly \cref{eq:ch18-gain-1d}. The one-dimensional Kalman update is therefore nothing but the inverse-variance-weighted average of two readings. (c) $\lambda^*=0.25/1.25=0.2$, $K=0.8$, $\hat x=0.2\cdot2+0.8\cdot3=2.8$~m and $\sigma^2=0.2$~m$^2$ ($\sigma=0.447$~m), the numbers of \cref{ex:ch18-1d}.
\end{solution}

\begin{solution}{exr:ch18-observability}
(a) $\mat{O}=\begin{psmallmatrix}\mat{H}\\ \mat{H}\mat{F}\end{psmallmatrix}=\begin{psmallmatrix}0&1\\0&1\end{psmallmatrix}$ has rank $1$: a velocity-only sensor never sees the position. The filter's position variance grows without bound --- it is the integral of the velocity error --- while the velocity variance settles, which is why dead reckoning drifts. (b) With $\mat{F}=\begin{psmallmatrix}1&\dt&\dt^2/2\\0&1&\dt\\0&0&1\end{psmallmatrix}$ and $\mat{H}=(1\;0\;0)$ the rows of $\mat{O}$ are $(1,0,0)$, $(1,\dt,\dt^2/2)$ and $(1,2\dt,2\dt^2)$, with determinant $\dt^3\neq0$: rank $3$, so position, velocity \emph{and} acceleration are all recoverable from a position sensor, as the CV case (rank $2$) already suggested. (c) $\mat{H}=(1\;0\;0\;0)$ on the state $(p_x,p_y,v_x,v_y)$ gives rank $2$; the observable subspace is spanned by $p_x$ and $v_x$, and $p_y,v_y$ are unobservable. \code{observability\_rank} returns $2$. In a simulation the $x$-block of $\mat{P}$ converges to the steady state of \cref{thm:ch18-steady} while the $y$-block follows the pure prediction recursion: $\operatorname{Var}(v_y)$ grows linearly in $t$ (like $qt$) and $\operatorname{Var}(p_y)$ like $qt^3/3$, so the reported uncertainty is honest --- the filter says it does not know $y$.
\end{solution}

\begin{solution}{exr:ch18-association}
Hint. The swap rate is set by the separation of the two predicted positions at the crossing measured in units of the innovation standard deviation. With $\sigma_p=1$~m and a $30^\circ$ crossing the two tracks stay within a couple of $\sigma_p$ of each other for several steps, and nearest-neighbour assignment is then close to a coin toss at each of those steps, so swaps are common; at $90^\circ$ the paths separate within one or two steps and swaps become rare, and halving $\sigma_p$ has the same effect as doubling the crossing angle's separation in $\sigma$ units. Count a swap by comparing the identity of each track's estimate with the truth after the intruders have separated by more than $10\sigma_p$, not at the crossing itself. The velocity-change test of part~(c) helps exactly in the ambiguous window and costs missed associations when the intruders really turn, so report both the swap rate and the number of dropped tracks.
\end{solution}

\begin{solution}{exr:ch18-coding}
Hint. With $\mat{H}=\mat{I}_4$ the innovation has $m=4$ components, so the mean NIS should sit near $4$, not $2$: the band of \cref{eq:ch18-nis-band} is $4\pm1.96\sqrt{8/N}$ over $N$ accepted updates, about $[3.6,4.4]$ for $N\approx180$, and the gate at $P_G=0.99$ is $\gamma=13.28$ from \cref{tab:ch18-chi2}. Compute the statistics only over accepted updates and only after the transient of the two-point initialisation, and remember that a dropped report is a predict-only step, not a zero innovation. The velocity channel makes the filter much better than in experiment~A, but it does not make finite differences respectable: the position-only filter is the fair baseline for part~(a), and the honest headline for part~(c) is the step-to-step change of the $3$~s extrapolation, which the raw reports make jump by metres and the filtered state by centimetres.
\end{solution}
